% C015 椭圆焦点三角形四心坐标性质——易错点
\begin{trapbox}
\begin{enumerate}
 \item \textbf{错：外心只在等腰三角形时位于 $y$ 轴。}

 \textbf{正：} 只要 $F_1=(-c,0)$、$F_2=(c,0)$ 是两个顶点，
 $F_1F_2$ 的垂直平分线就是 $y$ 轴；无论 $P$ 怎样移动，
 外心都在这条线上。\textbf{等腰条件只决定更多的对称性，不决定这个结论是否成立。}

 \item \textbf{错：四心坐标全是关于 $(x_0,y_0)$ 的线性变换。}

 \textbf{正：} $G,I$ 是固定比例的线性坐标变换；
 $H,O'$ 带有 $1/y_0$，一般是非线性变换。
 因此不能不加证明就声称垂心、外心的轨迹也是椭圆。

 \item \textbf{错：把内切圆半径 $r$ 当作内心纵坐标 $y_I$。}

 \textbf{正：} $r=c|y_0|/(a+c)>0$；
 若 $P$ 在 $x$ 轴下方，内心也在下方：$y_I=-r$。
 统一公式是 $\displaystyle y_I=cy_0/(a+c)$。

 \item \textbf{错：使用对边加权内心公式时把边长配错顶点。}

 \textbf{正：} $P$ 的对边是 $2c$，$F_1$ 的对边是 $PF_2$，
 $F_2$ 的对边是 $PF_1$。正确式子是
 \[
 I=\frac{2cP+PF_2\,F_1+PF_1\,F_2}{2(a+c)}.
 \]
 此处顶点记号表示坐标向量。

 \item \textbf{错：把 $PF_1$、$PF_2$ 的符号记反。}

 \textbf{正：} $P$ 越靠右，离左焦点越远：
 \[
 PF_1=a+ex_0,\qquad PF_2=a-ex_0.
 \]
 最可靠的现场推导是\textbf{和为 $2a$、平方差为 $4cx_0$}。

 \item \textbf{错：认为内心、顶角角平分线与底边交点是同一点。}

 \textbf{正：} 内心 $I$ 不在底边上；内切圆底边切点 $T$ 与角平分线落点 $D$ 也一般不是同一点。
 本题 $x_T=x_I=ex_0$，但 $x_D=e^2x_0$。

 \item \textbf{错：用斜率法时忽略竖直边。}

 \textbf{正：} 若 $x_0=c$，$PF_2$ 是竖直线，斜率不存在。
 改用方向向量点积 $\overrightarrow{F_1H}\cdot\overrightarrow{F_2P}=0$，
 就不需要除以 $x_0-c$。

 \item \textbf{错：一见分母 $y_0$ 就仍取 $y_0=0$。}

 \textbf{正：} $y_0=0$ 时椭圆上的 $P$ 只能是 $(\pm a,0)$，
 与两焦点共线，三角形退化，通常的内心、外心、垂心不再作为
 非退化三角形的四心存在。轨迹表达式须删去对应的端点。

 \item \textbf{错：认为 $P$ 在上半平面，外心必在上半平面。}

 \textbf{正：} 外心的纵坐标取决于 $OP^2-c^2$ 的符号：
 \[
 y_{O'}=\frac{OP^2-c^2}{2y_0}.
 \]
 当 $P$ 位于上半平面但 $OP<c$ 时，$O'$ 在下半平面；
 此时焦点三角形的顶角为钝角。

 \item \textbf{错：$x_I=ex_0$ 就意味着 $I$ 在 $P$ 到 $y$ 轴的水平线段上。}

 \textbf{正：} $|x_I|<|x_0|$（$x_0\ne0$），只能说明横坐标更靠近 $y$ 轴；
 实际上 $y_I=cy_0/(a+c)$，纵坐标也改变了。
 $I$ 通常\textbf{不与 $P$ 等高}，也通常不在 $OP$ 上。

 \item \textbf{错：将 $a=b$ 看成四心公式的普通圆特例。}

 \textbf{正：} $a=b$ 导致 $c=0$，两焦点重合，焦点三角形的底边消失。
 可以借鉴一般三角形的几何定理，但不能把退化圆当作非退化三角形。

 \item \textbf{错：等腰三角形的四心都重合。}

 \textbf{正：} 当 $P=(0,\pm b)$ 时四心都在 $y$ 轴上，
 但一般位于不同高度；只有三角形恰好是等边三角形时，四心才重合。
 本题等边条件是 $b=\sqrt3c$（等价于 $e=1/2$）。
\end{enumerate}
\end{trapbox}
